\par\SourceLeaf{512}%
\[
\begin{array}{c|}
\begin{array}{cc}\dfrac7{23\pdot73}&\dfrac{1680^2}{*}\end{array}\\[4pt]
\begin{array}{cc|}\dfrac{17}{23\pdot73}&\dfrac{136000}{*}\\[4pt]\dfrac{7\pdot73}{17}&\dfrac{511^2}{*}\end{array}\quad\text{aus diesen }\frac7{23}
\end{array}\quad\begin{gathered}\text{hieraus }73\\[3pt]\left[\text{und }\frac7{23},\ \frac7{17}\right]\end{gathered}
\]
\[
\begin{gathered}
\begin{array}{cc|}7\pdot7\pdot43&\dfrac{512001}{*}\\[4pt]17\pdot43&\dfrac{730^2}{*}\end{array}\quad
\frac{7\pdot7}{17}\ \text{hieraus und }\frac7{17}\left[\text{und }\frac7{23}\right]\\[3pt]
\text{wird }7,\ 17,\ 23,\ 43
\end{gathered}
\]\par\medskip\noindent\rule{.32\textwidth}{.3pt}
\begingroup\renewcommand{\arraystretch}{2.15}
\begin{longtable}{@{}llrr@{}}
$a,$ & $\frac{81}{8}$ & $\frac{81}{80}$ & \\
$b,$ & $\frac{41}{4}$ & $\frac{6561}{6560}$ & \\
$c,$ & $2$ & $\frac{1025}{1024}$ & \\
$d,$ & $7$ & $\frac{2401}{2400}$ & \\
$e,$ & $29$ & $\frac{1681^2}{1680\pdot 1682}$ & \\
$f,$ & $43$ & $\frac{512001}{512000}$ & \\
$g,$ & $23\pdot 73$ & $\frac{1680^2}{1679\pdot 1681}$ & \\
$h,$ & $17$ & $\frac{136000}{135999}$ & \\
$i,$ & $73$ & $\frac{511^2}{510\pdot 512}$ &  $\frac{730^2}{729\pdot 731}$\\
$k,$ & $13$ & $\frac{729^2}{728\pdot 730}$ & \\
$l,$ & $19$ & $\frac{512^2}{511\pdot 513}$ & \\
$m,$ & $47$ & $\frac{2116^2}{2115\pdot 2117}$ & \\
$n,$ & $61$ & $\frac{2500^2}{2499\pdot 2501}$ & \\
$o,$ & $31$ & $\frac{17081^2}{17080\pdot 17082}$ & \\
$p,$ & $11$ & $\frac{1024^2}{1023\pdot 1025}$ &  $\frac{2001^2}{2000\pdot 2002}$\\
$q,$ & $71$ & $\frac{10935^2}{10934\pdot 10936}$ & \\
$r,$ & $53$ & $\frac{788800}{788799}$ & \\
$s,$ & $67$ & $\frac{274700}{274699}$ &  $\frac{3751^2}{*}$\\
$t,$ & $37$ & $\frac{1000000}{999999}$ &  $\frac{1331^2}{1330\pdot 1332}$\\
$u,$ & $59$ & $\frac{3481^2}{3480\pdot 3481}$ & \\
$v,$ & $89$ & $\frac{4095^2}{4094\pdot 4096}$ & \\
$w,$ & $83$ & $\frac{6889^2}{6888\pdot 6890}$ & \\
$x,$ & $79$ & $\frac{3879^2}{3880\pdot 3882}$ & \\
$y,$ & $97$ & $\frac{13871^2}{13870\pdot 13872}$ &  $\frac{46656^2}{46655\pdot 46657}$\\
\end{longtable}\endgroup
[Die Anwendung dieser Brüche zur Bestimmung der Logarithmen der nebenstehenden kleinen Primzahlen mit Hilfe der noch wachsenden Potenzen von $\frac1x$ fortschreitenden sehr rasch konvergierenden Reihen für $\log\frac{x}{x-1}$ ergibt sich unmittelbar aus dem zu Anfang ausgeführten Beispiel.\par
Es lassen sich übrigens zur Bestimmung der Logarithmen der kleinsten Primzahlen $2$, $3$, $7$ noch vorteilhaftere Reihen aufstellen, wenn man diese Zahlen von \textsc{Gauss} auf geeignete Weise mit den von \textsc{Huyghens} (\textsc{Hugenii}, Opera varia, Lugduni $1724$ S. $457$) angegebenen verbindet:
\begingroup\renewcommand{\arraystretch}{1.15}
\begin{longtable}{@{}rl@{}}
$2$ & $1000$\\
 & $1024=2^{10}$\\
$3$ & $32805=3^8\pdot 5$\\
 & $32768=2^{15}$\\
$7$ & $2400=100\pdot 2^3\pdot 3$\\
 & $2401=7^4$\\
$11$ & $9800=100\pdot 2\pdot 7^2$\\
 & $9801=3^4\pdot 11^2$\\
$13$ & $123200=100\pdot 2^4\pdot 7\pdot 11$\\
 & $123201=3^6\pdot 13$\\
$17$ & $2600=100\pdot 2\pdot 13$\\
 & $2601=3^2\pdot 17^2$\\
$19$ & $28899=3^2\pdot 13^2\pdot 19$\\
 & $28900=100\pdot 17^2$\\
$23$ & $25920=10\pdot 2^5\pdot 3^4$\\
 & $25921=7^2\pdot 23^2$\\
$29$ & $613088=2^5\pdot 7^2\pdot 17\pdot 23$\\
 & $613089=3^6\pdot 29^2$\\
\end{longtable}\endgroup
